%% 04-the-monitoring-system.tex -- generated by tools/build_from_docx.py; edit freely

\hypertarget{the-monitoring-system}{%
\section{The Monitoring System}\label{the-monitoring-system}}

The Monitoring System is the capability through which an agent acquires
evidence about the objects it is registered to observe. It converts
telemetry, counters, events, state reports, logs and probe results
produced by the network and its management planes into typed, scoped,
quality-annotated observations carried to the Analytics Engine of the
same agent, or exposed to other agents through a declared observation
contract. Every proposition an agent later asserts, every decision it
reaches and every action it applies rests on what its Monitoring System
produced; the evidentiary quality of the loop is therefore fixed here
and cannot be recovered downstream.

One boundary governs the whole chapter. The Monitoring System produces
evidence and never conclusions. A measured value, a counter increment, a
state transition, a threshold crossing declared in the observation
contract and a log line are evidence. A score, a forecast, a
classification, an anomaly judgment and a root-cause hypothesis are not:
they are the output of analysis and belong to the Analytics Engine,
whatever function produced them and wherever that function resides. This
boundary is what makes the distinction between measurement and inference
auditable at the point where an agent's conclusions are challenged.

\hypertarget{role-and-scope-of-the-monitoring-system}{%
\subsection{Role and Scope of the Monitoring System}\label{role-and-scope-of-the-monitoring-system}}

\hypertarget{position-in-the-loop}{%
\subsubsection{Position in the loop}\label{position-in-the-loop}}

The Monitoring System occupies the first position of the MS-AE-DE-ACT
loop and holds two relationships with the rest of the framework. Toward
the network it is a consumer of management services, telemetry
interfaces and data products, bound to the sources its descriptor
declares and to no others. Toward the agent it is a producer of
observations, delivered over the local runtime binding where the
capabilities are co-located, or through the Data Continuum where they
are not. Observations are never carried on the service-based management
bus, which carries control only.

The Monitoring System also serves the Actuator. Where a verification
method declared by an action type reads the network back after an action
has been applied, that read is performed through the Monitoring System
of the same agent, against the same scope and the same quality rules as
any other observation. Verification evidence is therefore of the same
kind and the same provenance as the evidence on which the decision was
taken.

\hypertarget{observation-scope}{%
\subsubsection{Observation scope}\label{observation-scope}}

An agent observes the objects named by its observation scope and no
others. The scope is declared in the descriptor, resolved at admission
against the managed objects the assigned domain exposes, and held by the
platform as a versioned view. The Monitoring System runtime enforces it:
an observation whose subject lies outside the resolved scope is rejected
at the binding boundary rather than published, and the rejection is
recorded.

The scope view carries a version issued by the DMO of the assigned
domain. When the domain's inventory changes --- an object is added,
removed or re-parented --- the DMO issues a new version, and
observations published against a superseded version are marked as such.
This makes the difference between an object that was not observed and an
object that did not exist at the time a determinate one.

\hypertarget{capability-declarations}{%
\subsubsection{Capability declarations}\label{capability-declarations}}

The Monitoring System is declared at one of four levels. At
\texttt{full} the agent acquires and publishes its own observations
across the whole of its declared scope. At \texttt{partial} it acquires
a subset, and the subset shall be enumerated in the observation
interface; an unenumerated partial declaration is rejected at admission.
At \texttt{delegated} the agent consumes a feed produced by another
agent or by a platform service, and the delegation shall name both the
provider identifier and the interface profile. At \texttt{none} the
agent acquires nothing, which is permitted only where the declared
function demonstrably requires no observation.

An agent that declares a Monitoring System and no further capability is
a sensor. It may be registered and it may publish, but it shall not be
registered as a closed loop, and no authority over any managed object is
granted to it.

\hypertarget{measurement-kinds}{%
\subsection{Measurement Kinds}\label{measurement-kinds}}

Every metric an agent observes is of one of seven kinds, declared in the
type catalog and fixed for that metric across the platform. The kind
determines which aggregations are meaningful, how the runtime treats
discontinuities, and what the absence of a value means.

A \texttt{gauge} is an instantaneous value that may rise or fall. A
\texttt{cumulative\_\allowbreak{}counter} increases monotonically from an origin and
may reset; the runtime detects a reset as a decrease and reports the
affected interval as indeterminate rather than imputing the missing
increment. A \texttt{delta\_\allowbreak{}counter} reports the increment since the
previous report and requires no reset handling. An \texttt{event} is a
discrete occurrence with a timestamp and no duration. A \texttt{state}
is a value that persists until it changes and is reported on change. A
\texttt{log} is a text record with structured fields. A
\texttt{computed} metric is derived by the runtime from other metrics
under the rules of Section 4.4.

Permitted aggregations follow from the kind. A gauge admits mean,
minimum, maximum and percentiles but not summation across time. A
counter admits summation of increments but not averaging of absolute
values. An event admits counting. A state admits duration in state and
transition counts. A log admits counting and field extraction. Table 4-1
states the kinds, their definitions and the aggregations each permits.
Applying an aggregation the kind does not permit is a validity failure
at admission, not a runtime warning.

\begin{wptable}{4-1}

\begin{longtable}[]{@{}
  >{\raggedright\arraybackslash}p{(\columnwidth - 4\tabcolsep) * \real{0.2501}}
  >{\raggedright\arraybackslash}p{(\columnwidth - 4\tabcolsep) * \real{0.3750}}
  >{\raggedright\arraybackslash}p{(\columnwidth - 4\tabcolsep) * \real{0.3750}}@{}}
\caption{Measurement kinds and permitted aggregations}\tabularnewline
\toprule\noalign{}
\begin{minipage}[b]{\linewidth}\raggedright
\textbf{Measurement kind}
\end{minipage} & \begin{minipage}[b]{\linewidth}\raggedright
\textbf{Definition}
\end{minipage} & \begin{minipage}[b]{\linewidth}\raggedright
\textbf{Permitted aggregation}
\end{minipage} \\
\midrule\noalign{}
\endfirsthead
\toprule\noalign{}
\begin{minipage}[b]{\linewidth}\raggedright
\textbf{Measurement kind}
\end{minipage} & \begin{minipage}[b]{\linewidth}\raggedright
\textbf{Definition}
\end{minipage} & \begin{minipage}[b]{\linewidth}\raggedright
\textbf{Permitted aggregation}
\end{minipage} \\
\midrule\noalign{}
\endhead
\bottomrule\noalign{}
\endlastfoot
\textbf{gauge} & Instantaneous value that may rise or fall. & mean, min,
max, percentiles; not summation over time. \\
\textbf{cumulative\_\allowbreak{}counter} & Monotonically increasing
from an origin; may reset. & sum of increments; a reset is detected as a
decrease and the interval reported indeterminate. \\
\textbf{delta\_\allowbreak{}counter} & Increment since the previous
report. & sum of increments; no reset handling required. \\
\textbf{event} & Discrete occurrence with a timestamp and no duration. &
count. \\
\textbf{state} & Value persisting until it changes; reported on change.
& duration in state, transition count. \\
\textbf{log} & Text record with structured fields. & count, field
extraction. \\
\textbf{computed} & Derived by the runtime from other metrics (Section
4.4). & as permitted for the kind of the result. \\
\end{longtable}

\end{wptable}

\hypertarget{acquisition-windowing-and-quality}{%
\subsection{Acquisition, Windowing and Quality}\label{acquisition-windowing-and-quality}}

\hypertarget{acquisition-patterns}{%
\subsubsection{Acquisition patterns}\label{acquisition-patterns}}

A source is bound under one of six acquisition patterns, and the pattern
is declared per flow. A \texttt{periodic\_\allowbreak{}stream} delivers samples at a
declared interval. A \texttt{poll} is a read the Monitoring System
initiates on its own cadence. An \texttt{on\_\allowbreak{}change} source reports only
when a value changes. An \texttt{event\_\allowbreak{}notification} source reports
discrete occurrences as they arise. A \texttt{file} source delivers
records in batches at a declared or irregular cadence. An
\texttt{active\_\allowbreak{}probe} generates traffic and measures the response.

The pattern is not a deployment detail. It determines what completeness
means for the flow, and therefore what the platform may and may not say
about missing data (Table 4-2).

\begin{wptable}{4-2}

\begin{longtable}[]{@{}
  >{\raggedright\arraybackslash}p{(\columnwidth - 6\tabcolsep) * \real{0.2375}}
  >{\raggedright\arraybackslash}p{(\columnwidth - 6\tabcolsep) * \real{0.3125}}
  >{\raggedright\arraybackslash}p{(\columnwidth - 6\tabcolsep) * \real{0.2750}}
  >{\raggedright\arraybackslash}p{(\columnwidth - 6\tabcolsep) * \real{0.1751}}@{}}
\caption{Acquisition patterns and the determinability of
loss}\tabularnewline
\toprule\noalign{}
\begin{minipage}[b]{\linewidth}\raggedright
\textbf{Acquisition pattern}
\end{minipage} & \begin{minipage}[b]{\linewidth}\raggedright
\textbf{Definition}
\end{minipage} & \begin{minipage}[b]{\linewidth}\raggedright
\textbf{Basis of expectation}
\end{minipage} & \begin{minipage}[b]{\linewidth}\raggedright
\textbf{Loss}
\end{minipage} \\
\midrule\noalign{}
\endfirsthead
\toprule\noalign{}
\begin{minipage}[b]{\linewidth}\raggedright
\textbf{Acquisition pattern}
\end{minipage} & \begin{minipage}[b]{\linewidth}\raggedright
\textbf{Definition}
\end{minipage} & \begin{minipage}[b]{\linewidth}\raggedright
\textbf{Basis of expectation}
\end{minipage} & \begin{minipage}[b]{\linewidth}\raggedright
\textbf{Loss}
\end{minipage} \\
\midrule\noalign{}
\endhead
\bottomrule\noalign{}
\endlastfoot
\textbf{periodic\_\allowbreak{}stream} & Samples at a declared interval.
& Expected count from the interval. & Determinable. \\
\textbf{poll} & Read initiated by the MS on its own cadence. & Completed
reads against attempted reads. & Determinable. \\
\textbf{on\_\allowbreak{}change} & Reported only when the value changes.
& None. & Not determinable. \\
\textbf{event\_\allowbreak{}notification} & Discrete occurrences
reported as they arise. & None. & Not determinable. \\
\textbf{file} & Records delivered in batches. & Source-supplied sequence
identifiers, where present. & Determinable where sequenced. \\
\textbf{active\_\allowbreak{}probe} & Traffic generated and the response
measured. & Responses against probes emitted. & Determinable. \\
\end{longtable}

\end{wptable}

\hypertarget{windows}{%
\subsubsection{Windows}\label{windows}}

Observations are aggregated into tumbling windows aligned to UTC, of a
duration declared per flow. A window closes at its boundary and is
published once. A message arriving after its window has closed is
delivered marked late, carrying the window it belongs to; closed windows
are never republished and never silently revised. Consumers therefore
never see a value change beneath them, and a late arrival is visible as
a late arrival rather than as an inconsistency.

An observation is stale when its age exceeds the lesser of the staleness
threshold declared in the failure policy and the freshness requirement
of the flow that consumes it. Staleness is a property carried in the
envelope, not a reason to withhold publication.

\hypertarget{loss-and-coverage}{%
\subsubsection{Loss and coverage}\label{loss-and-coverage}}

Loss is defined per acquisition pattern, because only some patterns
carry the information from which loss can be determined. For a periodic
stream the runtime compares the samples received in a window against the
samples the declared interval implies. For a polled source it compares
completed reads against attempted reads. For a file source it uses the
sequence identifiers the source supplies, where it supplies them. For an
active probe it compares responses against probes emitted.

For an on-change source and for an event source no expected count
exists. Silence from such a source is indistinguishable, at the
Monitoring System, from a source that has stopped reporting. The runtime
shall therefore report loss for these patterns as not determinable, and
shall not report it as zero. A quality field that says nothing is safer
than one that asserts completeness the platform cannot establish.

Coverage is reported per metric rather than per flow, and is computed by
the runtime from the resolved scope: the declared cardinality, the
cardinality observed in the interval, and the enumerated identifiers of
the scope members from which nothing was received. Coverage is never
self-reported by contributed logic. For metrics acquired on change or by
event, coverage is established from source liveness --- the source's own
health signal or its declared heartbeat --- rather than from the arrival
of values, since for these patterns the arrival of no value is a
legitimate observation.

\hypertarget{quality-signals-and-de-duplication}{%
\subsubsection{Quality signals and de-duplication}\label{quality-signals-and-de-duplication}}

Every published observation carries loss as defined above, delivery
latency reported as the median, the ninety-ninth percentile and the
maximum over the window, and jitter computed as the difference between
the ninety-ninth percentile and the median of delivery latency. These
are properties of delivery, distinct from the measured quantity, and are
computed by the runtime.

Where the same quantity arrives from more than one source,
de-duplication uses the triple of managed object, metric and event time.
The order in which sources are declared in the flow establishes
precedence, so that the outcome of de-duplication is determined by the
descriptor rather than by arrival order.

A hard invariant governs all of the above and is not declarable: the
Monitoring System never fabricates a measurement. No interpolation, no
carry-forward and no synthetic fill is ever presented as measured. Gaps
are published as gaps.

\hypertarget{derivation-and-computed-metrics}{%
\subsection{Derivation and Computed Metrics}\label{derivation-and-computed-metrics}}

\hypertarget{where-computation-resides}{%
\subsubsection{Where computation resides}\label{where-computation-resides}}

Aggregation over a declared window and the computation of
catalog-defined metrics are performed by the Monitoring System runtime,
not by contributed logic. The consequence is that a metric name carries
one meaning across the platform: where two agents contributed by two
partners both publish a ninety-fifth percentile of a latency metric, the
two values were computed by the same code over windows of the same
definition, and are therefore comparable. A metric computed inside a
contributed function would carry that function's conventions with it,
and no consumer could establish what it meant.

Metric formulas of the kind standardized in 3GPP TS 28.554, and domain
formulas such as physical resource block utilization and blocking
probability, are held in the type catalog with their units and their
permitted aggregations, and are computed by the runtime from the catalog
entry.

\hypertarget{the-derivation-scheme}{%
\subsubsection{The derivation scheme}\label{the-derivation-scheme}}

Derivation is declared structurally and evaluated by the runtime. No
expression language is admitted, for the same reason none is admitted
for predicates: an expression language admits constructs whose cost and
whose behavior cannot be bounded at admission.

A derivation is one of two forms. An arithmetic derivation names an
operator --- \texttt{rate}, \texttt{delta}, \texttt{ratio},
\texttt{difference}, \texttt{sum} or \texttt{scale} --- and its
operands, each of which is a metric within the observation scope or a
literal. A roll-up derivation aggregates one metric across the objects
contained within a named object, under \texttt{sum\_\allowbreak{}over},
\texttt{mean\_\allowbreak{}over}, \texttt{min\_\allowbreak{}over}, \texttt{max\_\allowbreak{}over} or
\texttt{count\_\allowbreak{}over}, optionally restricted by a structured predicate of
the form defined for the descriptor, whose operands name object
attributes only. The operators are listed in Table 4-3. Derivations nest
to a depth of at most three.

\begin{wptable}{4-3}

\begin{longtable}[]{@{}
  >{\raggedright\arraybackslash}p{(\columnwidth - 4\tabcolsep) * \real{0.2125}}
  >{\raggedright\arraybackslash}p{(\columnwidth - 4\tabcolsep) * \real{0.1625}}
  >{\raggedright\arraybackslash}p{(\columnwidth - 4\tabcolsep) * \real{0.6250}}@{}}
\caption{Derivation operators}\tabularnewline
\toprule\noalign{}
\begin{minipage}[b]{\linewidth}\raggedright
\textbf{Operator}
\end{minipage} & \begin{minipage}[b]{\linewidth}\raggedright
\textbf{Form}
\end{minipage} & \begin{minipage}[b]{\linewidth}\raggedright
\textbf{Meaning}
\end{minipage} \\
\midrule\noalign{}
\endfirsthead
\toprule\noalign{}
\begin{minipage}[b]{\linewidth}\raggedright
\textbf{Operator}
\end{minipage} & \begin{minipage}[b]{\linewidth}\raggedright
\textbf{Form}
\end{minipage} & \begin{minipage}[b]{\linewidth}\raggedright
\textbf{Meaning}
\end{minipage} \\
\midrule\noalign{}
\endhead
\bottomrule\noalign{}
\endlastfoot
\textbf{rate} & Arithmetic & Change in an operand per unit time. \\
\textbf{delta} & Arithmetic & Difference between consecutive window
values of one operand. \\
\textbf{ratio} & Arithmetic & One operand divided by another. \\
\textbf{difference} & Arithmetic & One operand subtracted from
another. \\
\textbf{sum} & Arithmetic & Sum of operands. \\
\textbf{scale} & Arithmetic & Operand multiplied by a literal. \\
\textbf{sum\_\allowbreak{}over} & Roll-up & Sum of one metric across
contained objects. \\
\textbf{mean\_\allowbreak{}over} & Roll-up & Mean of one metric across
contained objects. \\
\textbf{min\_\allowbreak{}over} & Roll-up & Minimum across contained
objects. \\
\textbf{max\_\allowbreak{}over} & Roll-up & Maximum across contained
objects. \\
\textbf{count\_\allowbreak{}over} & Roll-up & Count of contained objects
satisfying a structured predicate. \\
\end{longtable}

\end{wptable}

A quantity that the scheme cannot express is not thereby unobtainable.
It is either work for the Analytics Engine, where inference is permitted
and where its provenance will be recorded as inference, or it is a
measurement the source itself shall supply --- as, for instance,
physical resource block fragmentation is supplied by the monitoring
application that computes it, rather than reconstructed from utilization
by the Monitoring System.

\hypertarget{validity-of-a-computed-metric}{%
\subsubsection{Validity of a computed metric}\label{validity-of-a-computed-metric}}

A computed metric is admitted only where four conditions hold together:
every metric it references lies within the agent's resolved observation
scope; the units of its operands and its result are those the type
catalog declares; each aggregation it applies is permitted for the kind
of the metric it is applied to; and every metric it references is
supplied by a source of a kind the flow declares. A derivation that
fails any of the four is rejected at admission rather than producing a
value at runtime whose meaning is undefined.

\hypertarget{sources-admitted-and-sources-excluded}{%
\subsection{Sources Admitted and Sources Excluded}\label{sources-admitted-and-sources-excluded}}

\hypertarget{analytics-outputs-are-not-monitoring-sources}{%
\subsubsection{Analytics outputs are not monitoring sources}\label{analytics-outputs-are-not-monitoring-sources}}

The output of an analytics function is not admissible as a Monitoring
System source. Network data analytics, computed statistics, forecasts
and predictions are inferences drawn from measurement, not measurement,
and the distinction does not depend on which function produced them: an
analytics output of the core network's analytics function and an
analytics output of a pipeline in the artificial-intelligence and
machine-learning domain are excluded on the same ground.

Such outputs reach an agent as Analytics Engine data products, or
through a delegated Analytics Engine, where the framework records them
as inference and where the agent's assurance record can show what rested
on measurement and what rested on a model. An agent that consumed a
prediction as though it were a measurement would present a conclusion
with the provenance of an observation, and nothing downstream could
recover the difference.

\hypertarget{generative-components-within-the-monitoring-system}{%
\subsubsection{Generative components within the Monitoring System}\label{generative-components-within-the-monitoring-system}}

A generative model may be used within the Monitoring System for two
purposes only. It may translate a query expressed in natural language
into a query against a declared source, and it may extract fields from
unstructured or semi-structured log records. Nothing else is permitted
within the capability.

Every value a generative component extracts shall be a verbatim span of
the source record, accompanied by a provenance pointer that identifies
the record and the position of the span within it. A value that is
paraphrased, normalized, corrected or inferred is not an extraction and
shall not be published as an observation. Diagnoses, summaries, severity
judgments and explanations of what a log indicates are the output of the
Analytics Engine, and a generative component within the Monitoring
System that produced them would be producing conclusions under the
provenance of evidence.

\hypertarget{the-global-monitoring-system}{%
\subsection{The Global Monitoring System}\label{the-global-monitoring-system}}

An agent may be assigned to the cross-domain management domain, where
its Monitoring System observes the management plane itself rather than
any single technological domain. Its sources are the registries the
cross-domain management function holds --- the liveness signals and
capability map of domain and capability registration, the service state
registry, and the outputs of service-level agreement monitoring ---
together with the reports domain management functions send northbound
under their delegations.

The global Monitoring System is distinct from the service-level
agreement monitor, which is a function of the cross-domain management
layer rather than a capability of an agent. The two stand in a defined
relation: the monitor may consume the global Monitoring System's
observations through its exposed contracts or as data products, and the
global Monitoring System may treat the monitor's outputs as a source in
the manner of any other management service. Neither subsumes the other.

The scope of the global Monitoring System is the management domain, and
the rules of this chapter apply to it without exception: its sources are
declared, its coverage is computed by the runtime against a resolved
scope, its loss is determinable only where the acquisition pattern
permits, and the substrate over which the underlying telemetry reaches
it is a property of the binding rather than of the capability.

\hypertarget{declared-failure-behavior}{%
\subsection{Declared Failure Behavior}\label{declared-failure-behavior}}

The Monitoring System's behavior under failure is declared rather than
left to the runtime's discretion, because the same degradation is
acceptable for one function and unacceptable for another. The failure
policy declares the four keys of Table 4-4.

It declares the response to source loss, as one of reduced sampling,
publication of a gap indicator only, degradation to cached values
explicitly marked as cached, or halting publication; where no response
is declared, publication of a gap indicator only applies. It declares
the staleness threshold beyond which an observation is marked stale. It
declares the response to partial scope loss, which shall in every case
publish a coverage indicator and shall never silently narrow the scope.
And it declares the cadence at which the Monitoring System publishes its
own health, independently of whether data is flowing, so that a silent
capability is distinguishable from an idle one.

\begin{wptable}{4-4}

\begin{longtable}[]{@{}
  >{\raggedright\arraybackslash}p{(\columnwidth - 6\tabcolsep) * \real{0.2625}}
  >{\raggedright\arraybackslash}p{(\columnwidth - 6\tabcolsep) * \real{0.1250}}
  >{\raggedright\arraybackslash}p{(\columnwidth - 6\tabcolsep) * \real{0.1000}}
  >{\raggedright\arraybackslash}p{(\columnwidth - 6\tabcolsep) * \real{0.5125}}@{}}
\caption{Monitoring System failure policy keys}\tabularnewline
\toprule\noalign{}
\begin{minipage}[b]{\linewidth}\raggedright
\textbf{Field in failure\_\allowbreak{}policy.\allowbreak{}MS}
\end{minipage} & \begin{minipage}[b]{\linewidth}\raggedright
\textbf{Type}
\end{minipage} & \begin{minipage}[b]{\linewidth}\raggedright
\textbf{Req.}
\end{minipage} & \begin{minipage}[b]{\linewidth}\raggedright
\textbf{Definition}
\end{minipage} \\
\midrule\noalign{}
\endfirsthead
\toprule\noalign{}
\begin{minipage}[b]{\linewidth}\raggedright
\textbf{Field in failure\_\allowbreak{}policy.\allowbreak{}MS}
\end{minipage} & \begin{minipage}[b]{\linewidth}\raggedright
\textbf{Type}
\end{minipage} & \begin{minipage}[b]{\linewidth}\raggedright
\textbf{Req.}
\end{minipage} & \begin{minipage}[b]{\linewidth}\raggedright
\textbf{Definition}
\end{minipage} \\
\midrule\noalign{}
\endhead
\bottomrule\noalign{}
\endlastfoot
\textbf{source\_\allowbreak{}loss} & enum & M &
reduced\_\allowbreak{}sampling,
gap\_\allowbreak{}indicator\_\allowbreak{}only,
degrade\_\allowbreak{}to\_\allowbreak{}cached or
halt\_\allowbreak{}publication.\allowbreak{} Default
gap\_\allowbreak{}indicator\_\allowbreak{}only.\allowbreak{} \\
\textbf{staleness\_\allowbreak{}threshold} & duration & M & Age beyond
which an observation is marked stale. \\
\textbf{partial\_\allowbreak{}scope\_\allowbreak{}loss} & enum & M &
Always publishes a coverage indicator; never silently narrows scope. \\
\textbf{health\_\allowbreak{}emission} & duration & M & Cadence of MS
self-health publication, independent of data flow. \\
\end{longtable}

\end{wptable}

\hypertarget{delegated-monitoring}{%
\subsection{Delegated Monitoring}\label{delegated-monitoring}}

Where an agent's Monitoring System is delegated, its evidence depends on
a provider it does not own and whose configuration may change without
reference to it. Three conditions convert that dependency into a checked
contract.

The delegation shall pin the provider's interface profile together with
a compatibility range. Admission shall verify that the provider's
declared observation scope covers the consumer's; a provider whose scope
is narrower is rejected rather than admitted with an unstated gap. And
the provider's health and coverage shall be forwarded to the consumer as
part of the observation binding itself, rather than published only to a
separate health channel that the consumer may or may not be reading.

A redeployment that narrows the provider's sampling or scope therefore
surfaces at the consumer as a binding fault, at the moment it occurs,
rather than as a slow and unattributable degradation in the quality of
the consumer's conclusions.

\hypertarget{subscription-and-adapter-boundaries}{%
\subsection{Subscription and Adapter Boundaries}\label{subscription-and-adapter-boundaries}}

Establishing a telemetry subscription is an act of reading, not an
action on the network. It is not subject to the action governance of the
Actuator, is not mediated, and consumes none of the agent's granted
control scope. It is nonetheless bounded: subscriptions are established
through source adapters the domain management function provides, those
adapters enforce read-only access to the underlying interface, and every
subscription an agent holds is removed when its generation ends.

This keeps the distinction between observing a network and changing one
structural rather than behavioral. An agent cannot acquire control of a
managed object through the act of subscribing to its telemetry, whatever
the underlying interface would otherwise permit.

\hypertarget{admission-consequences}{%
\subsection{Admission Consequences}\label{admission-consequences}}

Two admission behaviors specific to the Monitoring System are recorded
here because they determine what an agent may claim.

The latency class the platform publishes for an agent is the class its
placement can achieve, not the class its descriptor declares. Admission
computes an effective loop latency from the declared distribution of the
capabilities and the declared class of each transport hop between them,
and compares it with the declared class. Where the placement cannot
achieve the declared class, admission either rejects the agent or admits
it at the achievable class, which is then what composition and conflict
logic use.

A provisioning flag in a service request that indicates a domain
monitoring capability is not a declaration of a Monitoring System.
Scope, quality policy and failure policy are read from the descriptor
alone, and a service request carries no agent content by which a
monitoring role could be asserted. The flag continues to serve its own
purpose, which is to provision a capability in a domain, and carries no
agent semantics.

\hypertarget{scope-tags-across-trust-boundaries}{%
\subsection{Scope Tags Across Trust Boundaries}\label{scope-tags-across-trust-boundaries}}

Observations carry scope tags --- slice, tenant, region, site and
technology --- because conflict detection and routing depend on them.
Where an observation crosses a trust boundary, tag values that identify
a tenant or a slice are replaced by pseudonyms at the egress of the
trust domain. The mapping is held by the management layer and is not
resolvable by subscribing agents.

Conflict detection requires only that tags be comparable for equality,
which pseudonyms preserve. No arbitration capability is lost, and
tenancy is not disclosed to agents that have no need of it.

\hypertarget{the-observation-contract}{%
\subsection{The Observation Contract}\label{the-observation-contract}}

An observation contract is the declaration under which a Monitoring
System publishes. It names what is published, over which classes of
managed object, from which sources, at what freshness and under which
schema, and it is the object against which a consumer binds. A consumer
pins the contract identifier together with a compatibility range, so
that a change in what is published surfaces as a binding fault rather
than as a change in meaning that the consumer cannot see.

The contract is declared in the descriptor and resolved at admission.
Its fields are given in Table 4-5. Two constraints on it are structural
rather than conventional: the permitted binding classes of an
observation contract are the local runtime binding and the Data
Continuum and never the service-based management bus, and every metric
named in the contract resolves to a type catalog entry whose kind, unit
and permitted aggregations govern what the contract may declare for it.

\begin{wptable}{4-5}

\begin{longtable}[]{@{}
  >{\raggedright\arraybackslash}p{(\columnwidth - 6\tabcolsep) * \real{0.2875}}
  >{\raggedright\arraybackslash}p{(\columnwidth - 6\tabcolsep) * \real{0.1250}}
  >{\raggedright\arraybackslash}p{(\columnwidth - 6\tabcolsep) * \real{0.1625}}
  >{\raggedright\arraybackslash}p{(\columnwidth - 6\tabcolsep) * \real{0.4249}}@{}}
\caption{Observation contract fields}\tabularnewline
\toprule\noalign{}
\begin{minipage}[b]{\linewidth}\raggedright
\textbf{Field in envelope}
\end{minipage} & \begin{minipage}[b]{\linewidth}\raggedright
\textbf{Type}
\end{minipage} & \begin{minipage}[b]{\linewidth}\raggedright
\textbf{Req.}
\end{minipage} & \begin{minipage}[b]{\linewidth}\raggedright
\textbf{Definition}
\end{minipage} \\
\midrule\noalign{}
\endfirsthead
\toprule\noalign{}
\begin{minipage}[b]{\linewidth}\raggedright
\textbf{Field in envelope}
\end{minipage} & \begin{minipage}[b]{\linewidth}\raggedright
\textbf{Type}
\end{minipage} & \begin{minipage}[b]{\linewidth}\raggedright
\textbf{Req.}
\end{minipage} & \begin{minipage}[b]{\linewidth}\raggedright
\textbf{Definition}
\end{minipage} \\
\midrule\noalign{}
\endhead
\bottomrule\noalign{}
\endlastfoot
\textbf{observation\_\allowbreak{}id} & uuid & M & Unique identifier of
the observation, assigned by the MS runtime. \\
\textbf{agent\_\allowbreak{}instance\_\allowbreak{}id} & string & M &
Instance that produced the observation. \\
\textbf{contract\_\allowbreak{}id} & string & M & Observation contract
and version under which the observation is published. \\
\textbf{scope\_\allowbreak{}version} & string & M & Version of the
resolved scope view issued by the DMO; an observation against a
superseded version is marked. \\
\textbf{managed\_\allowbreak{}object} & string & M & Identifier of the
subject object, which lies within the resolved scope. \\
\textbf{metric} & string & M & Catalog identifier of the metric, with
its version. \\
\textbf{kind} & enum & M & Measurement kind from Table 4-1. \\
\textbf{unit} & string & M & Unit declared by the catalog entry for the
metric. \\
\textbf{value} & typed & C: not a gap & The measured or computed value.
Absent where the observation reports a gap. \\
\textbf{event\_\allowbreak{}time} & timestamp & M & Time to which the
value refers, on the normalized time base. \\
\textbf{window\_\allowbreak{}start} & timestamp & C: aggregated & Start
of the tumbling window, aligned to UTC. \\
\textbf{window\_\allowbreak{}duration} & duration & C: aggregated &
Declared duration of the window. \\
\textbf{ingest\_\allowbreak{}time} & timestamp & M & Time the MS runtime
received the underlying record. \\
\textbf{acquisition\_\allowbreak{}pattern} & enum & M & Pattern of the
flow that produced the observation (Table 4-2). \\
\textbf{derivation\_\allowbreak{}ref} & string & C: kind is computed &
Identifier of the declared derivation the runtime evaluated. \\
\textbf{scope\_\allowbreak{}tags} & object & M & Slice, tenant, region,
site and technology tags; tenant and slice values are pseudonymized
across a trust boundary. \\
\textbf{scope\_\allowbreak{}coverage} & object & M & Declared
cardinality, observed cardinality and the enumerated unobserved scope
members. Computed by the runtime. \\
\textbf{quality} & object & M & Loss (or
not\_\allowbreak{}determinable), delivery latency p50, p99 and maximum,
and jitter. \\
\textbf{stale} & bool & M & Whether the age exceeds the lesser of the
staleness threshold and the consuming flow\textquotesingle s freshness
requirement. \\
\textbf{late} & bool & M & Whether the record arrived after its window
had closed. \\
\textbf{gap} & bool & M & Whether the observation reports an absence of
measurement rather than a value. \\
\textbf{provenance} & object & C: generative extraction & Source record
identifier and the position of the verbatim span. \\
\textbf{source\_\allowbreak{}refs} & array & M & Sources contributing to
the value, in declared precedence order. \\
\end{longtable}

\end{wptable}

A contract is versioned independently of the descriptor that declares
it. Adding a metric or widening an object class is a compatible change;
removing a metric, narrowing an object class, changing a unit, changing
a window duration or changing the kind of a metric is not, and shall
increment the incompatible component of the version. This is what allows
a provider to evolve without silently invalidating the conclusions of
the agents that consume it.

\hypertarget{the-observation-envelope}{%
\subsection{The Observation Envelope}\label{the-observation-envelope}}

Every observation carries the envelope of Table 4-6. The envelope is
populated by the Monitoring System runtime and not by contributed logic,
and an observation reaching a binding boundary without a complete
envelope is rejected there rather than published.

\begin{wptable}{4-6}

\begin{longtable}[]{@{}
  >{\raggedright\arraybackslash}p{(\columnwidth - 6\tabcolsep) * \real{0.2875}}
  >{\raggedright\arraybackslash}p{(\columnwidth - 6\tabcolsep) * \real{0.1250}}
  >{\raggedright\arraybackslash}p{(\columnwidth - 6\tabcolsep) * \real{0.1625}}
  >{\raggedright\arraybackslash}p{(\columnwidth - 6\tabcolsep) * \real{0.4249}}@{}}
\caption{Observation envelope fields}\tabularnewline
\toprule\noalign{}
\begin{minipage}[b]{\linewidth}\raggedright
\textbf{Field in observation\_\allowbreak{}contract}
\end{minipage} & \begin{minipage}[b]{\linewidth}\raggedright
\textbf{Type}
\end{minipage} & \begin{minipage}[b]{\linewidth}\raggedright
\textbf{Req.}
\end{minipage} & \begin{minipage}[b]{\linewidth}\raggedright
\textbf{Definition}
\end{minipage} \\
\midrule\noalign{}
\endfirsthead
\toprule\noalign{}
\begin{minipage}[b]{\linewidth}\raggedright
\textbf{Field in observation\_\allowbreak{}contract}
\end{minipage} & \begin{minipage}[b]{\linewidth}\raggedright
\textbf{Type}
\end{minipage} & \begin{minipage}[b]{\linewidth}\raggedright
\textbf{Req.}
\end{minipage} & \begin{minipage}[b]{\linewidth}\raggedright
\textbf{Definition}
\end{minipage} \\
\midrule\noalign{}
\endhead
\bottomrule\noalign{}
\endlastfoot
\textbf{contract\_\allowbreak{}id} & string & M & Identifier of the
observation contract. \\
\textbf{version} & string & M & Version of the contract; consumers pin a
compatibility range. \\
\textbf{metrics{[}{]}} & array & M & Metrics published, each naming a
catalog entry and version. \\
\textbf{metrics{[}{]}.kind} & enum & M & Measurement kind, which shall
match the catalog entry. \\
\textbf{metrics{[}{]}.window} & duration & C: aggregated & Declared
window duration for the metric. \\
\textbf{metrics{[}{]}.aggregation} & enum & C: aggregated & Aggregation
applied, which shall be permitted for the kind. \\
\textbf{metrics{[}{]}.derivation} & object & C: kind is computed &
Structured derivation evaluated by the runtime (Table 4-3). \\
\textbf{object\_\allowbreak{}classes{[}{]}} & array & M & Classes of
managed object to which the metrics apply. \\
\textbf{sources{[}{]}} & array & M & Declared sources with their
acquisition pattern and interface, in precedence order. \\
\textbf{freshness} & duration & M & Maximum age at which a consumer may
rely on the observation. \\
\textbf{binding\_\allowbreak{}classes{[}{]}} & array & M & Permitted
bindings: LOCAL\_\allowbreak{}RUNTIME or
DATA\_\allowbreak{}CONTINUUM.\allowbreak{} The SMB is not permitted. \\
\textbf{schema\_\allowbreak{}ref} & string & M & Semantic schema of the
payload in the schema registry, with its version. \\
\end{longtable}

\end{wptable}

Table A-1 gives the descriptor location at which each field of Tables
4-4 and 4-6 is declared.

Three properties of the envelope carry most of its evidentiary weight.
The scope version and the coverage structure together establish what was
observed against what was in scope to observe, so that an absence is
attributable to a specific set of objects rather than reported as a
general shortfall. The separation of event time from ingest time keeps
the time at which something happened distinct from the time the platform
learned of it, which is what permits late arrivals to be marked rather
than silently reordered. And the gap flag, with the value absent rather
than defaulted, is the envelope-level expression of the invariant that
the Monitoring System never fabricates a measurement.

Where an observation is computed, the derivation reference identifies
the declared derivation the runtime evaluated, so that a consumer
challenging a computed value can establish exactly what was computed
from what. Where a value was extracted from a log by a generative
component, the provenance object identifies the source record and the
position of the verbatim span within it, and a consumer can retrieve the
original text.

\hypertarget{requirements-1}{%
\subsection{Requirements}\label{requirements-1}}

The requirements of this section apply to every agent that declares a
Monitoring System in any mode other than none, and to the platform that
admits and runs it.

\begin{requirement}{Req-T23-MS-01}{MUST}{Evidence, never conclusions}
\reqfield{Description / rationale}{%
The MS shall publish only measurement: values, counter increments,
events, states, declared threshold crossings and log records. A score,
forecast, classification, anomaly judgment or root-cause hypothesis
shall not be published as an observation, and admission shall reject an
observation contract whose schema carries model outputs, scores or
confidence fields.}
\reqfield{Dependency}{%
Req-T23-FW-02}
\reqfield{Traceability (Forward)}{%
\begin{itemize}
\tightlist
\item
  Checked by the validity rules against the declared observation schema.
\item
  Enforced at the binding boundary by the MS runtime.
\end{itemize}}
\reqfield{UNITY-6G components}{%
\begin{itemize}
\tightlist
\item
  Tuple descriptor
\item
  MS runtime
\item
  Admission (IDMO/DMO)
\end{itemize}}
\reqfield{Linked KPIs}{%
\begin{itemize}
\tightlist
\item
  KPI\#1: Deliver four distributed management entities (MS, AE, DE and
  ACT) deployable in each heterogeneous domain.
\item
  KPI\#36: Guarantee 99.99999\% service availability and reliability
  based on effective monitoring and LCM.
\end{itemize}}
\end{requirement}

\begin{requirement}{Req-T23-MS-02}{MUST}{Observation scope enforcement}
\reqfield{Description / rationale}{%
The MS shall acquire and publish only within the observation scope
resolved at admission and held by the platform as a DMO-versioned view.
An observation whose subject lies outside the resolved scope shall be
rejected at the binding boundary and the rejection recorded.
Observations published against a superseded scope version shall be
marked as such.}
\reqfield{Dependency}{%
Req-T23-MS-01}
\reqfield{Traceability (Forward)}{%
\begin{itemize}
\tightlist
\item
  Enforced by the MS runtime against the resolved scope.
\item
  Scope version issued by the DMO of the assigned domain.
\end{itemize}}
\reqfield{UNITY-6G components}{%
\begin{itemize}
\tightlist
\item
  MS runtime
\item
  DMO (inventory and scope view)
\end{itemize}}
\reqfield{Linked KPIs}{%
\begin{itemize}
\tightlist
\item
  KPI\#4: Real-time reconfiguration of IEEE 802.11 network domain in
  range smaller than 10 ms from taking the decision.
\item
  KPI\#36: Guarantee 99.99999\% service availability and reliability
  based on effective monitoring and LCM.
\item
  KPI\#1: Deliver four distributed management entities (MS, AE, DE and
  ACT) deployable in each heterogeneous domain.
\item
  KPI\#26: Increase the ratio of service management and LCM tasks
  resolved by local AE/DE components.
\item
  KPI\#33: Support more tenants via distributed AI \& DLT.
\end{itemize}}
\end{requirement}

\begin{requirement}{Req-T23-MS-03}{MUST}{No fabricated measurement}
\reqfield{Description / rationale}{%
The MS shall never present an interpolated, carried-forward or
synthetically filled value as measured. Gaps shall be published as gaps.
This invariant is not declarable and shall not be relaxed by any failure
policy.}
\reqfield{Dependency}{%
Req-T23-MS-01}
\reqfield{Traceability (Forward)}{%
\begin{itemize}
\tightlist
\item
  Enforced by the MS runtime; not expressible in the descriptor.
\end{itemize}}
\reqfield{UNITY-6G components}{%
\begin{itemize}
\tightlist
\item
  MS runtime
\end{itemize}}
\reqfield{Linked KPIs}{%
\begin{itemize}
\tightlist
\item
  KPI\#4: Real-time reconfiguration of IEEE 802.11 network domain in
  range smaller than 10 ms from taking the decision.
\item
  KPI\#5: Real-time reconfiguration of cellular network domain in the
  range smaller than 10 ms from taking the decision.
\item
  KPI\#36: Guarantee 99.99999\% service availability and reliability
  based on effective monitoring and LCM.
\end{itemize}}
\end{requirement}

\begin{requirement}{Req-T23-MS-04}{MUST}{Runtime aggregation and catalog metrics}
\reqfield{Description / rationale}{%
Window aggregation and the computation of catalog-defined metrics shall
be performed by the MS runtime and not by contributed logic, so that a
metric name denotes one computation across the platform. Catalog entries
shall carry the metric kind, its unit and its permitted aggregations,
and the runtime shall compute from the catalog entry.}
\reqfield{Dependency}{%
Req-T23-MS-02}
\reqfield{Traceability (Forward)}{%
\begin{itemize}
\tightlist
\item
  Computed by the MS runtime.
\item
  Catalog entry governs unit and permitted aggregation.
\end{itemize}}
\reqfield{UNITY-6G components}{%
\begin{itemize}
\tightlist
\item
  MS runtime
\item
  Type catalog (Annex B)
\end{itemize}}
\reqfield{Linked KPIs}{%
\begin{itemize}
\tightlist
\item
  KPI\#5: Real-time reconfiguration of cellular network domain in the
  range smaller than 10 ms from taking the decision.
\item
  KPI\#7: Cross-domain (cellular -- IEEE 802.11) end-to-end latency
  smaller than 20 ms.
\item
  KPI\#23: Reduction of communications overhead by a factor of 10 while
  maintaining the SLA requirements.
\item
  KPI\#8: Cross-domain jitter of max of 2 milliseconds, with less than
  10\% of the packets missing deadline.
\end{itemize}}
\end{requirement}

\begin{requirement}{Req-T23-MS-05}{MUST}{Structured derivation only}
\reqfield{Description / rationale}{%
A computed metric shall be declared as a structured derivation: an
arithmetic operator (rate, delta, ratio, difference, sum, scale) over
operands naming in-scope metrics or literals, or a roll-up
(sum\_\allowbreak{}over, mean\_\allowbreak{}over,
min\_\allowbreak{}over, max\_\allowbreak{}over,
count\_\allowbreak{}over) across contained objects, optionally
restricted by a structured predicate over object attributes. Nesting
shall not exceed depth three. No expression language shall be admitted.}
\reqfield{Dependency}{%
Req-T23-MS-04}
\reqfield{Traceability (Forward)}{%
\begin{itemize}
\tightlist
\item
  Checked by the validity rules at admission.
\item
  Evaluated by the MS runtime.
\end{itemize}}
\reqfield{UNITY-6G components}{%
\begin{itemize}
\tightlist
\item
  Tuple descriptor
\item
  MS runtime
\end{itemize}}
\reqfield{Linked KPIs}{%
\begin{itemize}
\tightlist
\item
  KPI\#5: Real-time reconfiguration of cellular network domain in the
  range smaller than 10 ms from taking the decision.
\item
  KPI\#7: Cross-domain (cellular -- IEEE 802.11) end-to-end latency
  smaller than 20 ms.
\end{itemize}}
\end{requirement}

\begin{requirement}{Req-T23-MS-06}{MUST}{Loss determinability by acquisition pattern}
\reqfield{Description / rationale}{%
Loss shall be computed from the expectation the declared acquisition
pattern establishes. Where the pattern establishes no expected count and
no sequence identifiers, in particular for on\_\allowbreak{}change and
event\_\allowbreak{}notification sources, loss shall be reported as not
determinable and shall not be reported as zero.}
\reqfield{Dependency}{%
Req-T23-MS-03}
\reqfield{Traceability (Forward)}{%
\begin{itemize}
\tightlist
\item
  Computed by the MS runtime per flow from the declared pattern.
\end{itemize}}
\reqfield{UNITY-6G components}{%
\begin{itemize}
\tightlist
\item
  MS runtime
\item
  Tuple descriptor
\end{itemize}}
\reqfield{Linked KPIs}{%
\begin{itemize}
\tightlist
\item
  KPI\#4: Real-time reconfiguration of IEEE 802.11 network domain in
  range smaller than 10 ms from taking the decision.
\item
  KPI\#8: Cross-domain jitter of max of 2 milliseconds, with less than
  10\% of the packets missing deadline.
\item
  KPI\#36: Guarantee 99.99999\% service availability and reliability
  based on effective monitoring and LCM.
\item
  KPI\#37: Guarantee 99.99999\% service continuity through AI-driven
  proactive management of cloud-native services.
\end{itemize}}
\end{requirement}

\begin{requirement}{Req-T23-MS-07}{MUST}{Platform-computed coverage}
\reqfield{Description / rationale}{%
Coverage shall be computed per metric by the MS runtime from the
resolved scope, reporting declared cardinality, observed cardinality for
the interval and the enumerated unobserved scope members, and shall
never be self-reported by contributed logic. For on\_\allowbreak{}change
and event\_\allowbreak{}notification metrics coverage shall be
established from source liveness rather than from value arrival.}
\reqfield{Dependency}{%
Req-T23-MS-02}
\reqfield{Traceability (Forward)}{%
\begin{itemize}
\tightlist
\item
  Computed at the MS runtime or Continuum ingress; envelopes lacking it
  are rejected.
\end{itemize}}
\reqfield{UNITY-6G components}{%
\begin{itemize}
\tightlist
\item
  MS runtime
\item
  Data Continuum ingress
\end{itemize}}
\reqfield{Linked KPIs}{%
\begin{itemize}
\tightlist
\item
  KPI\#4: Real-time reconfiguration of IEEE 802.11 network domain in
  range smaller than 10 ms from taking the decision.
\item
  KPI\#8: Cross-domain jitter of max of 2 milliseconds, with less than
  10\% of the packets missing deadline.
\item
  KPI\#33: Support more tenants via distributed AI \& DLT.
\item
  KPI\#36: Guarantee 99.99999\% service availability and reliability
  based on effective monitoring and LCM.
\end{itemize}}
\end{requirement}

\begin{requirement}{Req-T23-MS-08}{MUST}{Window discipline}
\reqfield{Description / rationale}{%
Observations shall be aggregated into UTC-aligned tumbling windows of a
duration declared per flow. A window shall be published once and never
republished or silently revised; a message arriving after its window has
closed shall be delivered marked late and carrying its window.}
\reqfield{Dependency}{%
Req-T23-MS-04}
\reqfield{Traceability (Forward)}{%
\begin{itemize}
\tightlist
\item
  Enforced by the MS runtime.
\end{itemize}}
\reqfield{UNITY-6G components}{%
\begin{itemize}
\tightlist
\item
  MS runtime
\end{itemize}}
\reqfield{Linked KPIs}{%
\begin{itemize}
\tightlist
\item
  KPI\#4: Real-time reconfiguration of IEEE 802.11 network domain in
  range smaller than 10 ms from taking the decision.
\item
  KPI\#8: Cross-domain jitter of max of 2 milliseconds, with less than
  10\% of the packets missing deadline.
\item
  KPI\#5: Real-time reconfiguration of cellular network domain in the
  range smaller than 10 ms from taking the decision.
\item
  KPI\#7: Cross-domain (cellular -- IEEE 802.11) end-to-end latency
  smaller than 20 ms.
\item
  KPI\#31: Reduce service recovery time below 180 s.
\item
  KPI\#34: Reduce Open RAN real-time control loops time scales below 10
  ms and 10x in NTN for real-time 6G.
\end{itemize}}
\end{requirement}

\begin{requirement}{Req-T23-MS-09}{MUST}{Analytics outputs excluded as sources}
\reqfield{Description / rationale}{%
The output of an analytics function, whether produced by a core network
analytics function, by an AI/ML domain pipeline or by any other
inference-producing function, shall not be admitted as an MS source.
Such outputs shall reach an agent as AE data products or through a
delegated AE, and shall be recorded as inference.}
\reqfield{Dependency}{%
Req-T23-MS-01}
\reqfield{Traceability (Forward)}{%
\begin{itemize}
\tightlist
\item
  Checked at admission against the declared source kinds.
\end{itemize}}
\reqfield{UNITY-6G components}{%
\begin{itemize}
\tightlist
\item
  Admission (IDMO/DMO)
\item
  Tuple descriptor
\end{itemize}}
\reqfield{Linked KPIs}{%
\begin{itemize}
\tightlist
\item
  KPI\#42: DT-validated AI algorithms that do not break the
  time-sensitive operation when deployed in 99\% of the outputs.
\item
  KPI\#26: Increase the ratio of service management and LCM tasks
  resolved by local AE/DE components.
\end{itemize}}
\end{requirement}

\begin{requirement}{Req-T23-MS-10}{MUST}{Confinement of generative components}
\reqfield{Description / rationale}{%
A generative model within the MS shall be used only to translate queries
against a declared source and to extract fields from log records. Every
extracted value shall be a verbatim span carrying a provenance pointer
to the record and the position of the span. Paraphrase, normalization,
correction, diagnosis and summary shall not be published as
observations.}
\reqfield{Dependency}{%
Req-T23-MS-01}
\reqfield{Traceability (Forward)}{%
\begin{itemize}
\tightlist
\item
  Checked at admission; enforced by the MS runtime at publication.
\end{itemize}}
\reqfield{UNITY-6G components}{%
\begin{itemize}
\tightlist
\item
  MS runtime
\item
  Admission (IDMO/DMO)
\end{itemize}}
\reqfield{Linked KPIs}{%
\begin{itemize}
\tightlist
\item
  KPI\#42: DT-validated AI algorithms that do not break the
  time-sensitive operation when deployed in 99\% of the outputs.
\item
  KPI\#26: Increase the ratio of service management and LCM tasks
  resolved by local AE/DE components.
\end{itemize}}
\end{requirement}

\begin{requirement}{Req-T23-MS-11}{MUST}{Declared failure behavior}
\reqfield{Description / rationale}{%
The MS failure policy shall declare the response to source loss
(reduced\_\allowbreak{}sampling,
gap\_\allowbreak{}indicator\_\allowbreak{}only,
degrade\_\allowbreak{}to\_\allowbreak{}cached,
halt\_\allowbreak{}publication; default
gap\_\allowbreak{}indicator\_\allowbreak{}only), a staleness threshold,
a response to partial scope loss that always publishes a coverage
indicator and never silently narrows scope, and a cadence of self-health
publication independent of data flow.}
\reqfield{Dependency}{%
Req-T23-MS-03}
\reqfield{Traceability (Forward)}{%
\begin{itemize}
\tightlist
\item
  Checked by the validity rules.
\item
  Enforced by the MS runtime.
\end{itemize}}
\reqfield{UNITY-6G components}{%
\begin{itemize}
\tightlist
\item
  Tuple descriptor
\item
  MS runtime
\end{itemize}}
\reqfield{Linked KPIs}{%
\begin{itemize}
\tightlist
\item
  KPI\#31: Reduce service recovery time below 180 s.
\item
  KPI\#33: Support more tenants via distributed AI \& DLT.
\item
  KPI\#4: Real-time reconfiguration of IEEE 802.11 network domain in
  range smaller than 10 ms from taking the decision.
\item
  KPI\#5: Real-time reconfiguration of cellular network domain in the
  range smaller than 10 ms from taking the decision.
\item
  KPI\#37: Guarantee 99.99999\% service continuity through AI-driven
  proactive management of cloud-native services.
\end{itemize}}
\end{requirement}

\begin{requirement}{Req-T23-MS-12}{MUST}{Delegated MS as a checked contract}
\reqfield{Description / rationale}{%
A delegated MS shall pin the provider interface profile and a
compatibility range; admission shall verify that the provider scope
covers the consumer scope; and provider health and coverage shall be
forwarded to the consumer within the observation binding. A provider
change leaving the pinned range shall raise a binding fault at the
consumer.}
\reqfield{Dependency}{%
Req-T23-FW-03}
\reqfield{Traceability (Forward)}{%
\begin{itemize}
\tightlist
\item
  Scope-subset check at admission; binding fault at runtime.
\end{itemize}}
\reqfield{UNITY-6G components}{%
\begin{itemize}
\tightlist
\item
  Admission (IDMO/DMO)
\item
  MS runtime
\item
  Provider agent
\end{itemize}}
\reqfield{Linked KPIs}{%
\begin{itemize}
\tightlist
\item
  KPI\#36: Guarantee 99.99999\% service availability and reliability
  based on effective monitoring and LCM.
\item
  KPI\#31: Reduce service recovery time below 180 s.
\item
  KPI\#26: Increase the ratio of service management and LCM tasks
  resolved by local AE/DE components.
\end{itemize}}
\end{requirement}

\begin{requirement}{Req-T23-MS-13}{MUST}{Subscriptions are reading}
\reqfield{Description / rationale}{%
Establishing a telemetry subscription shall not be treated as an action,
shall not be mediated and shall not consume granted control scope.
Subscriptions shall be established through DMO-provided source adapters
that enforce read-only access, and every subscription held by an agent
shall be removed when its generation ends.}
\reqfield{Dependency}{%
Req-T23-MS-02}
\reqfield{Traceability (Forward)}{%
\begin{itemize}
\tightlist
\item
  Enforced by the DMO source adapter and the lifecycle controller.
\end{itemize}}
\reqfield{UNITY-6G components}{%
\begin{itemize}
\tightlist
\item
  DMO (source adapters)
\item
  Lifecycle controller
\end{itemize}}
\reqfield{Linked KPIs}{%
\begin{itemize}
\tightlist
\item
  KPI\#26: Increase the ratio of service management and LCM tasks
  resolved by local AE/DE components.
\item
  KPI\#36: Guarantee 99.99999\% service availability and reliability
  based on effective monitoring and LCM.
\end{itemize}}
\end{requirement}

\begin{requirement}{Req-T23-MS-14}{MUST}{Global MS sources}
\reqfield{Description / rationale}{%
An MS assigned to the cross-domain management domain shall take as
sources the registries of domain and capability registration, the
service state registry, SLA monitoring outputs and northbound DMO
delegation reports. It shall remain distinct from the SLA monitor, which
may consume its output, and the rules of this chapter shall apply to it
without exception.}
\reqfield{Dependency}{%
Req-T23-MS-02}
\reqfield{Traceability (Forward)}{%
\begin{itemize}
\tightlist
\item
  Declared in the descriptor; checked at admission.
\end{itemize}}
\reqfield{UNITY-6G components}{%
\begin{itemize}
\tightlist
\item
  IDMO registries
\item
  MS runtime
\end{itemize}}
\reqfield{Linked KPIs}{%
\begin{itemize}
\tightlist
\item
  KPI\#23: Reduction of communications overhead by a factor of 10 while
  maintaining the SLA requirements.
\item
  KPI\#34: Reduce Open RAN real-time control loops time scales below 10
  ms and 10x in NTN for real-time 6G.
\end{itemize}}
\end{requirement}

\hypertarget{contribution-to-the-project-kpis}{%
\subsection{CONTRIBUTION TO THE PROJECT KPIs}\label{contribution-to-the-project-kpis}}

Table 4-7 states what the Monitoring System contributes to each KPI it
bears on, and the requirements through which the contribution is made.
Table 4-8 maps the requirements of this chapter onto those KPIs.

\begin{wptable}{4-7}

\begin{longtable}[]{@{}
  >{\raggedright\arraybackslash}p{(\columnwidth - 4\tabcolsep) * \real{0.1499}}
  >{\raggedright\arraybackslash}p{(\columnwidth - 4\tabcolsep) * \real{0.6375}}
  >{\raggedright\arraybackslash}p{(\columnwidth - 4\tabcolsep) * \real{0.2126}}@{}}
\caption{Contribution of the Monitoring System to the project
KPIs}\tabularnewline
\toprule\noalign{}
\begin{minipage}[b]{\linewidth}\raggedright
\textbf{KPI}
\end{minipage} & \begin{minipage}[b]{\linewidth}\raggedright
\textbf{Contribution of the Monitoring System}
\end{minipage} & \begin{minipage}[b]{\linewidth}\raggedright
\textbf{Requirements}
\end{minipage} \\
\midrule\noalign{}
\endfirsthead
\toprule\noalign{}
\begin{minipage}[b]{\linewidth}\raggedright
\textbf{KPI}
\end{minipage} & \begin{minipage}[b]{\linewidth}\raggedright
\textbf{Contribution of the Monitoring System}
\end{minipage} & \begin{minipage}[b]{\linewidth}\raggedright
\textbf{Requirements}
\end{minipage} \\
\midrule\noalign{}
\endhead
\bottomrule\noalign{}
\endlastfoot
\textbf{KPI\#1} & Supplies the MS of the four-entity tuple in every
heterogeneous domain, under one capability contract independent of
domain technology. & Req-T23-MS-01, Req-T23-MS-02 \\
\textbf{KPI\#4} & Bounds the evidence path for sub-10 ms IEEE 802.11
reconfiguration through local binding, declared windows and
runtime-computed quality. & Req-T23-MS-08, Req-T23-MS-11 \\
\textbf{KPI\#5} & Bounds the evidence path for sub-10 ms cellular
reconfiguration on the same terms. & Req-T23-MS-08, Req-T23-MS-11 \\
\textbf{KPI\#7} & Cross-domain latency is measured against a common time
base and a common window definition, so that an end-to-end figure means
the same in both domains. & Req-T23-MS-04, Req-T23-MS-08 \\
\textbf{KPI\#8} & Jitter is computed by the runtime as the difference
between the 99th percentile and the median of delivery latency, and
deadline misses remain countable because late messages are marked rather
than absorbed. & Req-T23-MS-04, Req-T23-MS-08 \\
\textbf{KPI\#23} & Aggregation over the declared window is performed by
the capability runtime at the source, so summaries travel rather than
raw high-rate telemetry, and a global MS reads the IDMO registries
rather than collecting again what the domains already hold. &
Req-T23-MS-04, Req-T23-MS-14 \\
\textbf{KPI\#42} & Evidence is separated from inference at the
capability boundary, and every generatively extracted value carries a
verbatim span and a provenance pointer, so no AI-derived value enters a
time-sensitive loop presented as measurement. & Req-T23-MS-09,
Req-T23-MS-10 \\
\textbf{KPI\#26} & A local MS with a local AE and DE keeps the evidence
path inside the domain, which is what permits a task to be resolved
locally. & Req-T23-MS-02, Req-T23-MS-12 \\
\textbf{KPI\#31} & Detection within the recovery budget rests on
declared windows, a declared staleness threshold and health emission
independent of data flow. & Req-T23-MS-08, Req-T23-MS-11 \\
\textbf{KPI\#33} & Scope tags crossing a trust boundary are
pseudonymized, so that multi-tenant observation discloses no tenancy
while remaining comparable for conflict detection. & Req-T23-MS-02,
Req-T23-MS-07 \\
\textbf{KPI\#34} & The effective latency class is computed from
placement rather than trusted from declaration, so that a sub-10 ms loop
is one the topology can achieve. & Req-T23-MS-08 \\
\textbf{KPI\#36} & Availability claims rest on coverage computed by the
platform against a resolved scope rather than asserted, on gaps
published as gaps, and on loss that cannot be determined being reported
as such rather than as zero. & Req-T23-MS-03, Req-T23-MS-06,
Req-T23-MS-07 \\
\textbf{KPI\#37} & Proactive management requires evidence whose
staleness and completeness are known, and both are envelope properties
computed by the runtime. & Req-T23-MS-06, Req-T23-MS-11 \\
\end{longtable}

\end{wptable}

\begin{wptable}{4-8}\footnotesize

\begin{longtable}[]{@{}
  >{\raggedright\arraybackslash}p{(\columnwidth - 26\tabcolsep) * \real{0.1813}}
  >{\raggedright\arraybackslash}p{(\columnwidth - 26\tabcolsep) * \real{0.0630}}
  >{\raggedright\arraybackslash}p{(\columnwidth - 26\tabcolsep) * \real{0.0630}}
  >{\raggedright\arraybackslash}p{(\columnwidth - 26\tabcolsep) * \real{0.0630}}
  >{\raggedright\arraybackslash}p{(\columnwidth - 26\tabcolsep) * \real{0.0630}}
  >{\raggedright\arraybackslash}p{(\columnwidth - 26\tabcolsep) * \real{0.0630}}
  >{\raggedright\arraybackslash}p{(\columnwidth - 26\tabcolsep) * \real{0.0630}}
  >{\raggedright\arraybackslash}p{(\columnwidth - 26\tabcolsep) * \real{0.0630}}
  >{\raggedright\arraybackslash}p{(\columnwidth - 26\tabcolsep) * \real{0.0630}}
  >{\raggedright\arraybackslash}p{(\columnwidth - 26\tabcolsep) * \real{0.0630}}
  >{\raggedright\arraybackslash}p{(\columnwidth - 26\tabcolsep) * \real{0.0630}}
  >{\raggedright\arraybackslash}p{(\columnwidth - 26\tabcolsep) * \real{0.0630}}
  >{\raggedright\arraybackslash}p{(\columnwidth - 26\tabcolsep) * \real{0.0630}}
  >{\raggedright\arraybackslash}p{(\columnwidth - 26\tabcolsep) * \real{0.0633}}@{}}
\caption{Mapping of Monitoring System requirements onto the project
KPIs}\tabularnewline
\toprule\noalign{}
\begin{minipage}[b]{\linewidth}\raggedright
\textbf{KPI\textbackslash{}}

\textbf{Requirement}
\end{minipage} & \begin{minipage}[b]{\linewidth}\raggedright
\textbf{\#1}
\end{minipage} & \begin{minipage}[b]{\linewidth}\raggedright
\textbf{\#4}
\end{minipage} & \begin{minipage}[b]{\linewidth}\raggedright
\textbf{\#5}
\end{minipage} & \begin{minipage}[b]{\linewidth}\raggedright
\textbf{\#7}
\end{minipage} & \begin{minipage}[b]{\linewidth}\raggedright
\textbf{\#8}
\end{minipage} & \begin{minipage}[b]{\linewidth}\raggedright
\textbf{\#23}
\end{minipage} & \begin{minipage}[b]{\linewidth}\raggedright
\textbf{\#26}
\end{minipage} & \begin{minipage}[b]{\linewidth}\raggedright
\textbf{\#31}
\end{minipage} & \begin{minipage}[b]{\linewidth}\raggedright
\textbf{\#33}
\end{minipage} & \begin{minipage}[b]{\linewidth}\raggedright
\textbf{\#34}
\end{minipage} & \begin{minipage}[b]{\linewidth}\raggedright
\textbf{\#36}
\end{minipage} & \begin{minipage}[b]{\linewidth}\raggedright
\textbf{\#37}
\end{minipage} & \begin{minipage}[b]{\linewidth}\raggedright
\textbf{\#42}
\end{minipage} \\
\midrule\noalign{}
\endfirsthead
\toprule\noalign{}
\begin{minipage}[b]{\linewidth}\raggedright
\textbf{KPI\textbackslash{}}

\textbf{Requirement}
\end{minipage} & \begin{minipage}[b]{\linewidth}\raggedright
\textbf{\#1}
\end{minipage} & \begin{minipage}[b]{\linewidth}\raggedright
\textbf{\#4}
\end{minipage} & \begin{minipage}[b]{\linewidth}\raggedright
\textbf{\#5}
\end{minipage} & \begin{minipage}[b]{\linewidth}\raggedright
\textbf{\#7}
\end{minipage} & \begin{minipage}[b]{\linewidth}\raggedright
\textbf{\#8}
\end{minipage} & \begin{minipage}[b]{\linewidth}\raggedright
\textbf{\#23}
\end{minipage} & \begin{minipage}[b]{\linewidth}\raggedright
\textbf{\#26}
\end{minipage} & \begin{minipage}[b]{\linewidth}\raggedright
\textbf{\#31}
\end{minipage} & \begin{minipage}[b]{\linewidth}\raggedright
\textbf{\#33}
\end{minipage} & \begin{minipage}[b]{\linewidth}\raggedright
\textbf{\#34}
\end{minipage} & \begin{minipage}[b]{\linewidth}\raggedright
\textbf{\#36}
\end{minipage} & \begin{minipage}[b]{\linewidth}\raggedright
\textbf{\#37}
\end{minipage} & \begin{minipage}[b]{\linewidth}\raggedright
\textbf{\#42}
\end{minipage} \\
\midrule\noalign{}
\endhead
\bottomrule\noalign{}
\endlastfoot
\textbf{Req-T23-MS-01} & \kpimark{} & & & & & & & & & & \kpimark{} &
& \\
\textbf{Req-T23-MS-02} & \kpimark{} & \kpimark{} & & & & & \kpimark{} &
& \kpimark{} & & \kpimark{} & & \\
\textbf{Req-T23-MS-03} & & \kpimark{} & \kpimark{} & & & & & & & &
\kpimark{} & & \\
\textbf{Req-T23-MS-04} & & & \kpimark{} & \kpimark{} & \kpimark{} &
\kpimark{} & & & & & & & \\
\textbf{Req-T23-MS-05} & & & \kpimark{} & \kpimark{} & & & & & & & &
& \\
\textbf{Req-T23-MS-06} & & \kpimark{} & & & \kpimark{} & & & & & &
\kpimark{} & \kpimark{} & \\
\textbf{Req-T23-MS-07} & & \kpimark{} & & & \kpimark{} & & & &
\kpimark{} & & \kpimark{} & & \\
\textbf{Req-T23-MS-08} & & \kpimark{} & \kpimark{} & \kpimark{} &
\kpimark{} & & & \kpimark{} & & \kpimark{} & & & \\
\textbf{Req-T23-MS-09} & & & & & & & \kpimark{} & & & & & &
\kpimark{} \\
\textbf{Req-T23-MS-10} & & & & & & & \kpimark{} & & & & & &
\kpimark{} \\
\textbf{Req-T23-MS-11} & & \kpimark{} & \kpimark{} & & & & & \kpimark{}
& \kpimark{} & & & \kpimark{} & \\
\textbf{Req-T23-MS-12} & & & & & & & \kpimark{} & \kpimark{} & & &
\kpimark{} & & \\
\textbf{Req-T23-MS-13} & & & & & & & \kpimark{} & & & & \kpimark{} &
& \\
\textbf{Req-T23-MS-14} & & & & & & \kpimark{} & & & & \kpimark{} & &
& \\
\end{longtable}

\end{wptable}
